% GENERATED FILE. Edit canonical lessons, then run npm run generate:books.
% canonical-source-hash: fnv1a64:2b62d052214bbb8f
% canonical-lessons: ES-C671-cercano, ES-C671-lejano, ES-C671-atajo, ES-C671-adelante, ES-C671-rodeo, ES-R671-repaso-664-671

\chapter{Cerca y lejos}
\label{ch:es-cerca-y-lejos}
\hlchaptermodality{671}

\begin{chapteropening}
\textbf{By the end of this chapter:} \emph{I can say how near or far a place is, and whether to take a short cut or the long way round.}

Complete the last lesson of chapter 671: I can say how near or far a place is, and whether to take a short cut or the long way round.
\end{chapteropening}

\section[cercano]{cercano — nearby, close}

\label{lesson:ES-C671-cercano}



\begin{quote}
Before the new one: say the Spanish for to cross, then the Spanish for to turn off.
\end{quote}

\subsection*{You'll want to know: cercano}
\textbf{cercano} — \textquotedblleft{}nearby, close\textquotedblright{}.

\begin{quote}
Busco una farmacia cercana.
\end{quote}

\begin{quote}
\emph{I am looking for a nearby chemist.}
\end{quote}

\emph{Un amigo cercano} — a close friend.

\begin{cousinweb}
\emph{cercano} is built on \emph{cerca} \textquotedblleft{}near\textquotedblright{} (EtymDB), which comes from Latin \emph{circā} \textquotedblleft{}around\textquotedblright{}.
\end{cousinweb}

\begin{grammarlens}[title={What you've built}]
One more word for telling, explaining and describing at B1.
\end{grammarlens}

\subsection*{Your turn}
\begin{itemize}
\raggedright
  \item \emph{Say it:} \emph{cercano}
  \item \emph{Say it:} \emph{cercano}, once more
  \item \emph{Say it:} a sentence of your own with \emph{cercano}
  \item \emph{Recall:} say the Spanish for to cross, then the Spanish for to turn off, then say \emph{cercano} again
\end{itemize}

\begin{tcolorbox}[breakable,colback=teal!4,colframe=teal!35!black,title={Before you move on}]
What does cercano mean? (\textquotedblleft{}Nearby, close\textquotedblright{}.) Say it once more.
\end{tcolorbox}
\section[lejano]{lejano — distant, far-off}

\label{lesson:ES-C671-lejano}



\begin{quote}
Before the new one: say the Spanish for to turn off, then the Spanish for nearby.
\end{quote}

\subsection*{You'll want to know: lejano}
\textbf{lejano} — \textquotedblleft{}distant, far-off\textquotedblright{}.

\begin{quote}
Viven en un pueblo lejano.
\end{quote}

\begin{quote}
\emph{They live in a far-off village.}
\end{quote}

The adverb is \emph{lejos}.

\begin{cousinweb}
\emph{lejano} is built on \emph{lejos} \textquotedblleft{}far\textquotedblright{} (EtymDB), which is usually traced to Latin \emph{laxius} \textquotedblleft{}more loosely, more widely\textquotedblright{}, from \emph{laxus} \textquotedblleft{}loose\textquotedblright{}.
\end{cousinweb}

\begin{grammarlens}[title={What you've built}]
One more word for telling, explaining and describing at B1.
\end{grammarlens}

\subsection*{Your turn}
\begin{itemize}
\raggedright
  \item \emph{Say it:} \emph{lejano}
  \item \emph{Say it:} \emph{lejano}, once more
  \item \emph{Say it:} a sentence of your own with \emph{lejano}
  \item \emph{Recall:} say the Spanish for to turn off, then the Spanish for nearby, then say \emph{lejano} again
\end{itemize}

\begin{tcolorbox}[breakable,colback=teal!4,colframe=teal!35!black,title={Before you move on}]
What does lejano mean? (\textquotedblleft{}Distant, far-off\textquotedblright{}.) Say it once more.
\end{tcolorbox}
\section[el atajo]{el atajo — the short cut}

\label{lesson:ES-C671-atajo}



\begin{quote}
Before the new one: say the Spanish for nearby, then the Spanish for distant.
\end{quote}

\subsection*{You'll want to know: el atajo}
\textbf{el atajo} — \textquotedblleft{}the short cut\textquotedblright{}.

\begin{quote}
Por este atajo llegamos antes.
\end{quote}

\begin{quote}
\emph{We get there sooner by this short cut.}
\end{quote}

\emph{Coger un atajo} — to take a short cut.

\begin{cousinweb}
\emph{atajo} is built on \emph{atajar} \textquotedblleft{}to cut short\textquotedblright{}, from \emph{tajar} \textquotedblleft{}to cut\textquotedblright{} (EtymDB), which comes from Late Latin \emph{taliāre} \textquotedblleft{}to cut\textquotedblright{}.
\end{cousinweb}

\begin{grammarlens}[title={What you've built}]
One more word for telling, explaining and describing at B1.
\end{grammarlens}

\subsection*{Your turn}
\begin{itemize}
\raggedright
  \item \emph{Say it:} \emph{el atajo}
  \item \emph{Say it:} \emph{el atajo}, once more
  \item \emph{Say it:} a sentence of your own with \emph{el atajo}
  \item \emph{Recall:} say the Spanish for nearby, then the Spanish for distant, then say \emph{el atajo} again
\end{itemize}

\begin{tcolorbox}[breakable,colback=teal!4,colframe=teal!35!black,title={Before you move on}]
What does el atajo mean? (\textquotedblleft{}The short cut\textquotedblright{}.) Say it once more.
\end{tcolorbox}
\section[adelante]{adelante — forward, ahead}

\label{lesson:ES-C671-adelante}



\begin{quote}
Before the new one: say the Spanish for distant, then the Spanish for the short cut.
\end{quote}

\subsection*{You'll want to know: adelante}
\textbf{adelante} — \textquotedblleft{}forward, ahead\textquotedblright{}.

\begin{quote}
Siga adelante hasta el cruce.
\end{quote}

\begin{quote}
\emph{Keep going to the junction.}
\end{quote}

\emph{¡Adelante!} also means come in! or go ahead!

\begin{cousinweb}
\emph{adelante} comes from \emph{delante} \textquotedblleft{}in front\textquotedblright{} and Old Spanish \emph{denante} (EtymDB), a fusion of Latin \emph{de}, \emph{in} and \emph{ante} \textquotedblleft{}before\textquotedblright{}.
\end{cousinweb}

\begin{grammarlens}[title={What you've built}]
One more word for telling, explaining and describing at B1.
\end{grammarlens}

\subsection*{Your turn}
\begin{itemize}
\raggedright
  \item \emph{Say it:} \emph{adelante}
  \item \emph{Say it:} \emph{adelante}, once more
  \item \emph{Say it:} a sentence of your own with \emph{adelante}
  \item \emph{Recall:} say the Spanish for distant, then the Spanish for the short cut, then say \emph{adelante} again
\end{itemize}

\begin{tcolorbox}[breakable,colback=teal!4,colframe=teal!35!black,title={Before you move on}]
What does adelante mean? (\textquotedblleft{}Forward, ahead\textquotedblright{}.) Say it once more.
\end{tcolorbox}
\section[el rodeo]{el rodeo — the detour, the long way round; the rodeo}

\label{lesson:ES-C671-rodeo}



\begin{quote}
Before the new one: say the Spanish for the short cut, then the Spanish for forward.
\end{quote}

\subsection*{You'll want to know: el rodeo}
\textbf{el rodeo} — \textquotedblleft{}the detour, the long way round; the rodeo\textquotedblright{}.

\begin{quote}
Si vamos por el puente, damos un rodeo.
\end{quote}

\begin{quote}
\emph{If we go over the bridge, we go the long way round.}
\end{quote}

\emph{Sin rodeos} — without beating about the bush.

\begin{cousinweb}
\emph{rodeo} is built on \emph{rodear} \textquotedblleft{}to go round\textquotedblright{}, from \emph{rueda} \textquotedblleft{}wheel\textquotedblright{}, which is inherited from Latin \emph{rota}.
\end{cousinweb}

\begin{grammarlens}[title={What you've built}]
One more word for telling, explaining and describing at B1.
\end{grammarlens}

\subsection*{Your turn}
\begin{itemize}
\raggedright
  \item \emph{Say it:} \emph{el rodeo}
  \item \emph{Say it:} \emph{el rodeo}, once more
  \item \emph{Say it:} a sentence of your own with \emph{el rodeo}
  \item \emph{Recall:} say the Spanish for the short cut, then the Spanish for forward, then say \emph{el rodeo} again
\end{itemize}

\begin{tcolorbox}[breakable,colback=teal!4,colframe=teal!35!black,title={Before you move on}]
What does el rodeo mean? (\textquotedblleft{}The detour, the long way round; the rodeo\textquotedblright{}.) Say it once more.
\end{tcolorbox}
\section[(dialogue)]{Repaso de los capítulos 664–671}

\label{lesson:ES-R671-repaso-664-671}



\begin{quote}
Say the Spanish for the sheet and the balcony.
\end{quote}

\subsection*{Your turn}
Say these aloud:

\begin{itemize}
\raggedright
  \item the sheet, the balcony, the view, the socket, the tap
  \item the complaint, the minor damage, to compensate, the manager, the owner
  \item the sleeping bag, to camp, the sleeping mat, the caravan, the torch
  \item the flat, the proprietor, the host, the tenant, to rent
  \item to undo, to pack up, the scales, the excess, to drag
  \item the street map, to find your way, the location, to point to, located
  \item to show, the instruction, to turn, to cross, to turn off
  \item nearby, distant, the short cut, forward, the detour
\end{itemize}

\begin{tcolorbox}[breakable,colback=teal!4,colframe=teal!35!black,title={Before you move on}]
The sheet? (\textbf{la sábana}.) The balcony? (\textbf{el balcón}.) The view? (\textbf{la vista}.) The socket? (\textbf{el enchufe}.) The tap? (\textbf{el grifo}.) The complaint? (\textbf{la queja}.) The minor damage? (\textbf{el desperfecto}.) To compensate? (\textbf{compensar}.) The manager? (\textbf{el gerente}.) The owner? (\textbf{el dueño}.) The sleeping bag? (\textbf{el saco de dormir}.) To camp? (\textbf{acampar}.) The sleeping mat? (\textbf{la esterilla}.) The caravan? (\textbf{la caravana}.) The torch? (\textbf{la linterna}.) The flat? (\textbf{el apartamento}.) The proprietor? (\textbf{el propietario}.) The host? (\textbf{el anfitrión}.) The tenant? (\textbf{el inquilino}.) To rent? (\textbf{alquilar}.) To undo? (\textbf{deshacer}.) To pack up? (\textbf{empaquetar}.) The scales? (\textbf{la báscula}.) The excess? (\textbf{el exceso}.) To drag? (\textbf{arrastrar}.) The street map? (\textbf{el plano}.) To find your way? (\textbf{orientarse}.) The location? (\textbf{la ubicación}.) To point to? (\textbf{señalar}.) Located? (\textbf{situado}.) To show? (\textbf{indicar}.) The instruction? (\textbf{la indicación}.) To turn? (\textbf{doblar}.) To cross? (\textbf{atravesar}.) To turn off? (\textbf{desviarse}.) Nearby? (\textbf{cercano}.) Distant? (\textbf{lejano}.) The short cut? (\textbf{el atajo}.) Forward? (\textbf{adelante}.) The detour? (\textbf{el rodeo}.)
\end{tcolorbox}
